\subsection*{用一个数描述整块区域的变化}
已知每点去哪之后，可以问：面积为一的方形，变换后占多大地方？两列向量围成的平行四边形就是它的像。若两列为 $(a,c)$ 与 $(b,d)$，计算有向面积会得到 $ad-bc$。这个数总结了变换对面积的一项整体作用。

剪切使上边向旁边滑动，但底与垂直高度不变，所以面积不变。水平反射也不改变面积大小，却翻转了两根基箭头的排列朝向。为了保存这个区别，我们让面积倍率带符号；负号表示定向翻转，不是实际土地具有负面积。

\subsection*{零值意味着信息丢失}
让两列逐渐变得平行，平行四边形的高度消失，面积归零。此时不同输入压到同一条线上，许多位置无法区分。行列式为零由此预告了信息丢失，也解释了为什么无法通过逆变换找回所有输入。

在矩阵实验中把第二列设为第一列的两倍。先预测面积，再检查显示值；随后轻微偏转它，观察面积恢复。图形可以帮助理解压扁，但“严格等于零”仍依赖精确数值或符号关系，不能仅凭图形很窄来判断。

复合的倍率应该相乘：第一步把每块面积扩大两倍，第二步再扩大三倍，总倍率就是六。定向也遵守乘法规则，两次反射恢复原朝向。这给出了行列式乘法公式的意义；下面的代数展开检查它在一般二维坐标下成立。面积、可逆性与零空间将在下一节汇合。
